
\documentclass[11pt,a4paper]{article}
\usepackage[T1]{fontenc}
\usepackage{lmodern}
\usepackage{amsmath,amssymb,amsfonts,mathtools}
\usepackage{graphicx,booktabs,longtable,array,tabularx,adjustbox}
\usepackage{geometry,microtype,fancyhdr,lastpage,etoolbox,enumitem}
\usepackage{newunicodechar,textcomp}
\usepackage{hyperref}
\hypersetup{hidelinks,pdfcreator={XeLaTeX / QM1 academic edition}}
\geometry{left=20mm,right=20mm,top=22mm,bottom=23mm}
\setcounter{secnumdepth}{2}
\setlength{\parindent}{0pt}
\setlength{\parskip}{4pt plus 1pt}
\setlength{\emergencystretch}{2em}
\setlength{\tabcolsep}{3pt}
\renewcommand{\arraystretch}{1.16}
\setlist{nosep,topsep=3pt,itemsep=2pt,leftmargin=7mm}
\AtBeginEnvironment{longtable}{\footnotesize\setlength{\tabcolsep}{2.2pt}}
\newcommand{\tightlist}{\setlength{\itemsep}{0pt}\setlength{\parskip}{0pt}}
\providecommand{\Cov}{\operatorname{Cov}}
\providecommand{\Var}{\operatorname{Var}}
\providecommand{\Corr}{\operatorname{Corr}}
\providecommand{\MAD}{\operatorname{MAD}}
\providecommand{\E}{\mathbb{E}}
\providecommand{\R}{\mathbb{R}}
\newunicodechar{—}{\textemdash}
\newunicodechar{–}{\textendash}
\newunicodechar{€}{\texteuro}
\newunicodechar{≤}{\ensuremath{\leq}}
\newunicodechar{≥}{\ensuremath{\geq}}
\newunicodechar{×}{\ensuremath{\times}}
\newunicodechar{÷}{\ensuremath{\div}}
\newunicodechar{Σ}{\ensuremath{\Sigma}}
\newunicodechar{σ}{\ensuremath{\sigma}}
\newunicodechar{μ}{\ensuremath{\mu}}
\newunicodechar{Δ}{\ensuremath{\Delta}}
\newunicodechar{∑}{\ensuremath{\sum}}
\newunicodechar{≈}{\ensuremath{\approx}}
\newunicodechar{√}{\ensuremath{\surd}}
\newunicodechar{≠}{\ensuremath{\ne}}
\newunicodechar{→}{\ensuremath{\to}}
\newunicodechar{∞}{\ensuremath{\infty}}
\newunicodechar{∈}{\ensuremath{\in}}
\newunicodechar{−}{\ensuremath{-}}
\newunicodechar{²}{\ensuremath{^{2}}}
\newunicodechar{³}{\ensuremath{^{3}}}
\newunicodechar{₀}{\ensuremath{_{0}}}
\newunicodechar{₁}{\ensuremath{_{1}}}
\newunicodechar{₂}{\ensuremath{_{2}}}
\newunicodechar{₃}{\ensuremath{_{3}}}
\newunicodechar{₄}{\ensuremath{_{4}}}
\newunicodechar{ₜ}{\ensuremath{_{t}}}
\newunicodechar{ₚ}{\ensuremath{_{p}}}
\newunicodechar{ₙ}{\ensuremath{_{n}}}
\newunicodechar{ₓ}{\ensuremath{_{x}}}
\newunicodechar{ᵢ}{\ensuremath{_{i}}}
\newunicodechar{ŷ}{\ensuremath{\hat y}}
\newunicodechar{Ŷ}{\ensuremath{\widehat Y}}
\newunicodechar{ₘ}{\ensuremath{_{m}}}
\newunicodechar{ₑ}{\ensuremath{_{e}}}
\newunicodechar{ₐ}{\ensuremath{_{a}}}
\newunicodechar{ₖ}{\ensuremath{_{k}}}
\newunicodechar{ₛ}{\ensuremath{_{s}}}
\newunicodechar{ₗ}{\ensuremath{_{l}}}
\newunicodechar{ₕ}{\ensuremath{_{h}}}
\newunicodechar{ₒ}{\ensuremath{_{o}}}
\newunicodechar{₊}{\ensuremath{_{+}}}
\newunicodechar{₋}{\ensuremath{_{-}}}
\newunicodechar{₌}{\ensuremath{_{=}}}
\newunicodechar{₍}{\ensuremath{_{(}}}
\newunicodechar{₎}{\ensuremath{_{)}}}
\newunicodechar{⁻}{\ensuremath{^{-}}}
\newunicodechar{⁺}{\ensuremath{^{+}}}
\newunicodechar{⁰}{\ensuremath{^{0}}}
\newunicodechar{¹}{\ensuremath{^{1}}}
\newunicodechar{⁴}{\ensuremath{^{4}}}
\newunicodechar{⁵}{\ensuremath{^{5}}}
\newunicodechar{⁶}{\ensuremath{^{6}}}
\newunicodechar{⁷}{\ensuremath{^{7}}}
\newunicodechar{⁸}{\ensuremath{^{8}}}
\newunicodechar{⁹}{\ensuremath{^{9}}}
\makeatletter
\renewcommand{\section}{\@startsection{section}{1}{0pt}{2.0ex plus .7ex minus .3ex}{1.0ex}{\large\bfseries}}
\renewcommand{\subsection}{\@startsection{subsection}{2}{0pt}{1.8ex plus .5ex minus .2ex}{.7ex}{\normalsize\bfseries}}
\renewcommand{\subsubsection}{\@startsection{subsubsection}{3}{0pt}{1.3ex plus .4ex minus .2ex}{.5ex}{\normalsize\bfseries}}
\makeatother
\pagestyle{fancy}
\fancyhf{}
\renewcommand{\headrulewidth}{0.35pt}
\renewcommand{\footrulewidth}{0pt}
\fancyhead[L]{\small\textsc{KEDGE Business School --- QM1}}
\fancyhead[R]{\small\textsc{2026--2027}}
\fancyfoot[L]{\small Badr ABOUFARASSE, PhD}
\fancyfoot[R]{\small\thepage\ /\ \pageref{LastPage}}
\setlength{\headheight}{14pt}
\setkeys{Gin}{keepaspectratio}
\begin{document}

\begin{center}
{\small\textsc{KEDGE Business School \quad EBP1 \quad Quantitative Methods 1}}\\[6pt]
{\Large\bfseries Worksheet 4 — Bivariate distributions and regression}\\[6pt]
{\normalsize\bfseries Fully worked solutions}\\[6pt]
{\small Badr ABOUFARASSE, PhD \quad --- \quad 2026--2027}
\end{center}
\vspace{2mm}\hrule\vspace{6mm}

\hypertarget{core-applications}{%
\section{Core applications}\label{core-applications}}

Master the essential definitions and calculations from the course.

\emph{Core application --- 15 min}

\hypertarget{exercise-3-mayer-line}{%
\subsection{Exercise 3 -- Mayer line}\label{exercise-3-mayer-line}}

For {\((1,7),(2,10),(3,10),(4,14),(5,14)\)}, form two sub-clouds, calculate mean points {\(A\)} and {\(B\)} and determine the Mayer line.

\hypertarget{fully-worked-solution}{%
\subsubsection{Fully worked solution}\label{fully-worked-solution}}

{\[A=\left(\frac{1+2+3}{3},\frac{7+10+10}{3}\right)=(2,9),\]} {\[B=\left(\frac{4+5}{2},\frac{14+14}{2}\right)=(4.5,14).\]} {\[a=\frac{14-9}{4.5-2}=2,
\qquad b=5.\]} {\[\boxed{y=2x+5}.\]}

\hypertarget{auditable-construction-of-the-two-centroids}{%
\subsubsection{Auditable construction of the two centroids}\label{auditable-construction-of-the-two-centroids}}

\begin{longtable}[]{@{}lllll@{}}
\toprule\noalign{}
Group & x & y & Mean x & Mean y \\
\midrule\noalign{}
\endhead
\bottomrule\noalign{}
\endlastfoot
A & 1, 2, 3 & 7, 10, 10 & 6/3=2 & 27/3=9 \\
B & 4, 5 & 14, 14 & 9/2=4.5 & 28/2=14 \\
\end{longtable}

\[a=(14-9)/(4.5-2)=2,\quad b=9-2\times2=5.\]

Checks: 2×2+5=9 and 2×4.5+5=14. The Mayer line is not necessarily the least-squares line.

\hypertarget{assessment-and-exam-practice}{%
\section{Assessment and exam practice}\label{assessment-and-exam-practice}}

Practise independently with assessment and final-exam style questions.

\emph{Assessment / exam practice --- 35 min}

\hypertarget{exercise-1-complete-linear-regression}{%
\subsection{Exercise 1 -- Complete linear regression}\label{exercise-1-complete-linear-regression}}

Consider {\((1,3),(2,5),(3,5),(4,8),(5,9)\)}.

\begin{enumerate}
\item
  Draw the scatter plot and describe direction, form and apparent strength.
\item
  Calculate {\(\bar x\)}, {\(\bar y\)}, {\(V(x)\)}, {\(V(y)\)}, {\(\Cov(x,y)\)} and correlation {\(r\)}.
\item
  Estimate the least-squares line {\(\hat y=ax+b\)}.
\item
  Calculate fitted values and residuals {\(e_i=y_i-\hat y_i\)}.
\item
  Calculate {\(SSE\)}, {\(SST\)} and {\(R^2\)} and interpret {\(R^2\)}.
\end{enumerate}

\hypertarget{fully-worked-solution-1}{%
\subsubsection{Fully worked solution}\label{fully-worked-solution-1}}

{\[\bar x=3,\qquad \bar y=6,\qquad
V(x)=2,\qquad V(y)=4.8.\]} {\[\Cov(x,y)=3,
\qquad
r=\frac{3}{\sqrt{2\times4.8}}\approx0.96825.\]} {\[a=\frac{3}{2}=1.5,
\qquad
b=6-1.5(3)=1.5.\]} {\[\boxed{\hat
y=1.5x+1.5}.\]}

\begin{longtable}[]{@{}llllll@{}}
\toprule\noalign{}
{\(x_i\)} & 1 & 2 & 3 & 4 & 5 \\
\midrule\noalign{}
\endhead
\bottomrule\noalign{}
\endlastfoot
{\(y_i\)} & 3 & 5 & 5 & 8 & 9 \\
{\(\hat y_i\)} & 3 & 4.5 & 6 & 7.5 & 9 \\
{\(e_i\)} & 0 & 0.5 & -1 & 0.5 & 0 \\
\end{longtable}

{\[SSE=1.5,
\qquad
SST=24,\]} {\[R^2=1-\frac{1.5}{24}=0.9375.\]} About 93.75\% of the observed variation in {\(Y\)} around its mean is accounted for by the fitted linear model.

Full computation table, including residuals and squared residuals:

\begin{longtable}[]{@{}lllllllll@{}}
\toprule\noalign{}
xᵢ & yᵢ & xᵢ² & yᵢ² & xᵢyᵢ & ŷᵢ & eᵢ & eᵢ² & \\
\midrule\noalign{}
\endhead
\bottomrule\noalign{}
\endlastfoot
1 & 3 & 1 & 9 & 3 & 3 & 0 & 0 & \\
2 & 5 & 4 & 25 & 10 & 4.5 & 0.5 & 0.25 & \\
3 & 5 & 9 & 25 & 15 & 6 & -1 & 1 & \\
4 & 8 & 16 & 64 & 32 & 7.5 & 0.5 & 0.25 & \\
5 & 9 & 25 & 81 & 45 & 9 & 0 & 0 & \\
Σ & 15 & 30 & 55 & 204 & 105 & 30 & 0 & 1.5 \\
\end{longtable}

\[\bar x=15/5=3,\quad\bar y=30/5=6,\quad V(x)=55/5-3^2=2,\quad V(y)=204/5-6^2=4.8.\]

\[\operatorname{Cov}(x,y)=105/5-3(6)=3,\quad a=3/2=1.5,\quad b=6-1.5(3)=1.5.\]

\[SSE=0+.25+1+.25+0=1.5,\quad SST=24,\quad R^2=1-1.5/24=.9375.\]

\includegraphics[width=0.85\textwidth,height=\textheight]{assets/chapter-4-solutions-en-chart-01.pdf}

\emph{Nuage de points / Scatter and OLS line}

\emph{Assessment / exam practice --- 15 min}

\hypertarget{exercise-4-two-stage-forecast}{%
\subsection{Exercise 4 -- Two-stage forecast}\label{exercise-4-two-stage-forecast}}

\begin{longtable}[]{@{}lllllllll@{}}
\toprule\noalign{}
Month & 1 & 2 & 3 & 4 & 5 & 6 & 7 & 8 \\
\midrule\noalign{}
\endhead
\bottomrule\noalign{}
\endlastfoot
Advertising {\(X\)} & 12 & 15 & 14 & 18 & 20 & 22 & 25 & 27 \\
Sales {\(Y\)} & 48 & 52 & 50 & 59 & 63 & 68 & 73 & 78 \\
\end{longtable}

Given {\(\widehat Y=2.064X+21.899\)} and {\(\widehat X_t=2.155t+9.429\)}, forecast month-10 advertising and sales, and identify where extrapolation occurs.

\hypertarget{fully-worked-solution-2}{%
\subsubsection{Fully worked solution}\label{fully-worked-solution-2}}

{\[\widehat
X_{10}=2.155(10)+9.429\approx30.98.\]} This exceeds the observed advertising maximum 27, so the first step extrapolates the time trend beyond month 8. {\[\widehat
Y_{10}=2.064(30.98)+21.899\approx85.84.\]} The final sales forecast is about 85.84 thousand euros.

\hypertarget{complete-two-stage-calculation}{%
\subsubsection{Complete two-stage calculation}\label{complete-two-stage-calculation}}

\begin{longtable}[]{@{}lllllllll@{}}
\toprule\noalign{}
t & 1 & 2 & 3 & 4 & 5 & 6 & 7 & 8 \\
\midrule\noalign{}
\endhead
\bottomrule\noalign{}
\endlastfoot
X & 12 & 15 & 14 & 18 & 20 & 22 & 25 & 27 \\
Y & 48 & 52 & 50 & 59 & 63 & 68 & 73 & 78 \\
\end{longtable}

\[\widehat X_{10}=2.155(10)+9.429=21.55+9.429=30.979.\]

\[\widehat Y_{10}=2.064(30.979)+21.899=63.940656+21.899=85.839656.\]

Month 10 is outside the observed range 1--8, and predicted advertising 30.979 exceeds the observed maximum 27: this is a two-stage extrapolation requiring caution.

\emph{Assessment / exam practice --- 45 min}

\hypertarget{exercise-5-two-successive-regressions-exam-style}{%
\subsection{Exercise 5 --- Two successive regressions --- exam style}\label{exercise-5-two-successive-regressions-exam-style}}

For six periods, observe period t, advertising x and sales y.

\begin{longtable}[]{@{}lllllll@{}}
\toprule\noalign{}
t & 1 & 2 & 3 & 4 & 5 & 6 \\
\midrule\noalign{}
\endhead
\bottomrule\noalign{}
\endlastfoot
x & 4 & 5 & 7 & 8 & 10 & 11 \\
y & 10 & 12 & 15 & 18 & 22 & 24 \\
\end{longtable}

\begin{enumerate}
\tightlist
\item
  Build the complete t, x, y, t², x², y², tx, xy table and sums.
\item
  Compute means, variances, covariance, r and the regression of x on t.
\item
  Repeat for y on x.
\item
  Forecast x then y for t=8; discuss extrapolation.
\end{enumerate}

\hypertarget{fully-worked-solution-3}{%
\subsubsection{Fully worked solution}\label{fully-worked-solution-3}}

\begin{longtable}[]{@{}llllllll@{}}
\toprule\noalign{}
t & x & y & t² & x² & y² & t×x & x×y \\
\midrule\noalign{}
\endhead
\bottomrule\noalign{}
\endlastfoot
1 & 4 & 10 & 1 & 16 & 100 & 4 & 40 \\
2 & 5 & 12 & 4 & 25 & 144 & 10 & 60 \\
3 & 7 & 15 & 9 & 49 & 225 & 21 & 105 \\
4 & 8 & 18 & 16 & 64 & 324 & 32 & 144 \\
5 & 10 & 22 & 25 & 100 & 484 & 50 & 220 \\
6 & 11 & 24 & 36 & 121 & 576 & 66 & 264 \\
Σt=21 & 45 & 101 & 91 & 375 & 1853 & 183 & 833 \\
\end{longtable}

Regression of x on t.

\[\bar t=21/6=3.5,\ \bar x=45/6=7.5,\ V(t)=91/6-3.5^2=2.916667,\ V(x)=375/6-7.5^2=6.25.\]

\[\begin{gathered}
\operatorname{Cov}(t,x)=183/6-3.5(7.5)=4.25,\\r_{tx}=\frac{4.25}{\sqrt{2.916667\times6.25}}=.995418.
\end{gathered}\]

\[a_1=4.25/2.916667=1.457143,\quad b_1=7.5-a_1(3.5)=2.4,\quad\widehat x_t=1.457143t+2.4.\]

Regression of y on x.

\[\begin{gathered}
\bar y=101/6=16.833333,\\V(y)=1853/6-(101/6)^2=25.472222,\\\operatorname{Cov}(x,y)=833/6-(45/6)(101/6)=12.583333.
\end{gathered}\]

\[\begin{gathered}
r_{xy}=\frac{12.583333}{\sqrt{6.25\times25.472222}}=.997292,\\a_2=12.583333/6.25=2.013333,\\b_2=16.833333-a_2(7.5)=1.733333.
\end{gathered}\]

\[\begin{gathered}
\widehat y=2.013333x+1.733333,\\\widehat x_8=14.057143,\\\widehat y_8=2.013333(14.057143)+1.733333=30.035048.
\end{gathered}\]

t=8 is beyond the observed 1--6 range. Both stages extrapolate; neither establishes causation.

\hypertarget{optional-advanced-practice}{%
\section{Optional advanced practice}\label{optional-advanced-practice}}

Explore additional interpretations and professional applications.

\emph{Advanced practice --- 20 min}

\hypertarget{exercise-2-interpolation-extrapolation-causality-and-influence}{%
\subsection{Exercise 2 -- Interpolation, extrapolation, causality and influence}\label{exercise-2-interpolation-extrapolation-causality-and-influence}}

Using Exercise 1, forecast {\(Y\)} at {\(X=3.5\)} and {\(X=10\)}, classify each forecast, add point {\((10,30)\)} and recalculate the regression line, then explain why high correlation does not establish causality.

\hypertarget{fully-worked-solution-4}{%
\subsubsection{Fully worked solution}\label{fully-worked-solution-4}}

{\[\hat y(3.5)=1.5(3.5)+1.5=6.75.\]} Because {\(3.5\in[1,5]\)}, this is interpolation. {\[\hat y(10)=16.5.\]} Since 10 lies far outside the observed range, this is risky extrapolation.

After adding {\((10,30)\)}, the new line is approximately {\[\hat
y=3.04918x-2.70492.\]} The slope rises from 1.5 to about 3.05, showing that the new observation is highly influential. High correlation remains an association and cannot by itself establish causality.

After adding (10,30), N=6, Σx=25, Σy=60, Σx²=155, Σy²=1104, Σxy=405.

\[V(x)=155/6-(25/6)^2=8.472222,\quad\operatorname{Cov}(x,y)=405/6-(25/6)(60/6)=25.833333.\]

\[a=25.833333/8.472222=3.049180,\quad b=10-3.049180(25/6)=-2.704918.\]

The slope changes from 1.50 to 3.05: the additional point is influential.

\emph{Advanced practice --- 20 min}

\hypertarget{extra-1-regression-diagnosis-and-responsible-decision}{%
\subsection{Extra 1 --- Regression diagnosis and responsible decision}\label{extra-1-regression-diagnosis-and-responsible-decision}}

A model linking store traffic to temperature reports r = 0.78 and R² = 0.61. Its residual plot is U-shaped, one summer weekend has a very large residual, and the manager wants to forecast a temperature 12°C above the observed maximum.

\begin{enumerate}
\tightlist
\item
  Interpret r and R² correctly.
\item
  Diagnose the residual pattern.
\item
  Discuss the unusual weekend and the proposed forecast.
\item
  Recommend the next analytical steps.
\end{enumerate}

\textbf{Method hint:}

Separate strength, functional form, unusual observations, extrapolation and causality.

\hypertarget{fully-worked-solution-5}{%
\subsubsection{Fully worked solution}\label{fully-worked-solution-5}}

\textbf{Correlation.} r = 0.78 indicates a fairly strong positive linear association in the observed sample. It does not establish causality.

\textbf{R².} About 61\% of the sample variation in traffic around its mean is explained by the fitted linear relationship. It does not mean that temperature ``causes 61\%'' of traffic.

\textbf{Residual diagnosis.} A U-shape is systematic structure, indicating that a straight line misses curvature. Predictions may be biased at low and high temperatures.

\textbf{Unusual weekend.} Check data quality and contextual factors such as a festival, promotion or closure. Report sensitivity with and without the observation; do not delete it mechanically.

\textbf{Extrapolation.} Forecasting 12°C beyond the maximum is unsupported and particularly unsafe because curvature is already visible.

\textbf{Next steps.} Add a quadratic term or use a flexible curve, validate on held-out periods, include calendar and promotional controls, inspect influence statistics, and present prediction uncertainty.

\end{document}
